The validity of the Agent-Based Model (ABM) relied heavily on the empirical grounding of its initialization parameters. Before training the Deep Reinforcement Learning (DRL) agent, the real-world data gathered from the Municipality of Bacolod—comprising legislative documents and key informant interviews—were analyzed to define the state space, constraints, and reward structures of the simulation environment. This section detailed how qualitative field data was translated into quantitative simulation parameters.

\subsection{Legislative Constraints and Financial Bounds}
Based on the simulation parameters defined in the barangay configuration, the interview transcripts referenced in Appendix \ref{ch:appendix_b}, and the financial data provided by the Municipal Environment and Natural Resources Office (MENRO), this section detailed the legislative and financial constraints that defined the solution space for the Agent-Based Model. The simulation environment was strictly bounded by the fiscal realities of the Municipality of Bacolod. Unlike theoretical models where resources might be infinite, this study incorporated hard constraints derived from Appropriation Ordinance No. 2024-01 (refer to Appendix \ref{appendix:appro_ordinance}), ensuring that the Deep Reinforcement Learning (DRL) agent operated within the actual legal and financial framework of the Local Government Unit (LGU).

At the macro-environmental level, the MENRO acted as the primary resource allocator. The simulation enforced a strict Annual Budget Cap of \textpeso 1,500,000, a figure calibrated directly from the MENRO interview detailed in Appendix \ref{appendix:menro_interview}. This created a quarterly resource constraint of \textpeso 375,000. Under the ``Status Quo'' policy, these funds were typically distributed based on population density; however, the simulation revealed that a purely demographic distribution diluted the impact in high-density areas while over-funding low-density areas. This confirmed the ``institutional failure'' described by \cite{Ibanez2022} and \cite{Santos2025}, who argued that systemic funding deficits fundamentally weakened the implementation of RA 9003 at the municipal level.

Furthermore, operational constraints were dictated by the regional economic reality. With a Daily Minimum Wage of \textpeso 400 based on Region X rates, the cost of human enforcement was significantly high. A single enforcer patrolling for a quarter consumed approximately \textpeso 24,000 to \textpeso 30,000. This financial reality created a ``Personnel Trap'' where the budget was consumed entirely by salaries, leaving zero room for Information, Education, and Communication (IEC) campaigns. This validated the observation by \cite{Nishimura2022} regarding the operational burden on LGUs, where the high cost of technical and human resources created a bottleneck that prevented holistic program implementation.

\begin{table}[htbp]
    \centering
    \caption{Barangay Demographic and Financial Initialization Parameters}
    \label{tab:barangay_demographics}
    \small 
    \begin{tabular}{l r r c c c c}
        \toprule
        & & {Annual} & {Self-Reported} & \multicolumn{3}{c}{{Income Profile (\%)}} \\
        \cmidrule(lr){5-7}
        {Barangay} & {Households} & {Budget (\textpeso)} & {Compliance} & {Low} & {Mid} & {High} \\
        \midrule
        Brgy Poblacion    & 1,534 & 200,000 & 2\%  & 70 & 25 & 5 \\
        Brgy Liangan East & 608   & 30,000  & 65\% & 40 & 40 & 20 \\
        Brgy Ezperanza    & 574   & 90,000  & 20\% & 20 & 50 & 30 \\
        Brgy Binuni       & 507   & 126,370 & 95\% & 50 & 30 & 20 \\
        Brgy Demologan    & 463   & 21,000  & 60\% & 80 & 15 & 5 \\
        Brgy Babalaya     & 171   & 15,000  & 100\% & 80 & 10 & 10 \\
        Brgy Mati         & 165   & 80,000  & 70\% & 90 & 5  & 5 \\
        \bottomrule
    \end{tabular}
    \vspace{1ex}
    
    \raggedright
    \footnotesize
    \textit{Note: Income profiles derived from Appendix B interviews. Low Income agents have higher price sensitivity ($\gamma$) to fines.}
\end{table}

The heterogeneity of the seven barangays created a complex optimization landscape where each operated under unique financial bounds and legislative priorities. Barangay Poblacion represented an urban resource trap.  Despite having the highest local budget of \textpeso 200,000 (refer to Appendix \ref{appendix:appro_ord_poblacion}), it suffered from the lowest per-capita spending power due to its massive population of 1,534 households. Defined as the ``Police State'' profile in Appendix \ref{appendix:poblacion_interview}, Poblacion was legislatively bound to allocate 60\% of its budget to enforcement because the high density of commerce and residence necessitated a heavy presence of Barangay Tanods to maintain basic order. The simulation demonstrated that despite the high absolute budget, this 60\% enforcement lock left insufficient funds for IEC or incentives, resulting in a rigid compliance ceiling of approximately 2\%. This aligned with findings by \cite{Collado2024}, who noted that while strict enforcement policies like ``No Segregation, No Collection'' existed, they were often unsustainable without complementary behavioral interventions.

Conversely, Barangay Binuni represented the simulation’s ``Best Case Scenario'' or wealthy anomaly. With a budget of \textpeso 126,370 for only 507 households (refer to Appendix \ref{appendix:appro_ord_binuni}), it enjoyed a high per-capita budget. This fiscal surplus allowed Binuni to adopt a balanced ``Rich \& Capable'' profile, as noted in Appendix \ref{appendix:binuni_interview}, splitting funds evenly between enforcement and incentives. This financial flexibility allowed Binuni to overcome high behavioral friction, making it the consistent top performer in the Status Quo analysis.

Barangay Liangan East faced a critical mismatch that highlighted the risks of incentive-heavy programs without sufficient funding. It attempted an ambitious program with a meager budget of \textpeso 30,000 for 608 households (refer to Appendix \ref{appendix:appro_ord_liangan-east}). Modeled after the ``Eco-brick'' project mentioned in Appendix \ref{appendix:liangan_interview}, Liangan East was legislatively biased toward incentives at 65\%, leaving only 25\% for enforcement. This created a fragile system where the simulation showed that without LGU augmentation, the \textpeso 30,000 was quickly exhausted by incentive redemptions, leading to rapid program collapse. This validated the ``False Spike'' behavioral profile, where compliance rose briefly but crashed when funds ran out. This phenomenon mirrored the conclusions of \cite{Camarillo2021}, who emphasized that while incentive programs were attractive, their long-term effectiveness was often volatile and dependent on continuous financial support.

In contrast, Barangay Ezperanza represented the middle-class standard. With \textpeso 90,000 for 574 households (refer to Appendix \ref{appendix:appro_ord_ezperanza}), it served as the median control group. Following a ``Consistent'' profile from Appendix \ref{appendix:ezperanza_interview}, it maintained a 50\% enforcement and 30\% incentive split, allowing it to maintain a baseline compliance of roughly 10-15\% without external intervention.

Barangay Mati acted as a small-scale surplus outlier with a high budget of \textpeso 80,000 for a tiny population of 165 households (refer to Appendix \ref{appendix:appro_ord_mati}). Being 90\% agricultural as described in Appendix \ref{appendix:mati_interview}, the legislative focus was on incentives rather than strict policing. The high budget-to-household ratio meant Mati was highly responsive to financial injections; however, the simulation assigned a high decay rate to this profile, indicating that while farmers participated for rewards, the habit formation was volatile.

Finally, Barangays Babalaya and Demologan represented the poverty trap of the simulation, operating with budgets of \textpeso 15,000 (refer to Appendix \ref{appendix:appro_ord_babalaya}) and \textpeso 21,000 (refer to Appendix \ref{appendix:appro_ord_demologan}) respectively. They were bound by the ``Survivor'' profile detailed in Appendices \ref{appendix:babalaya_interview} and \ref{appendix:demologan_interview}. With such limited funds, 85\% to 90\% of the budget was locked into enforcement, not as a strategic choice but a survival necessity. \textpeso 15,000 was barely enough to pay the honoraria of two Tanods for a quarter. Consequently, their IEC and incentive allocation was effectively zero. The simulation confirmed that under the Status Quo, these barangays were mathematically incapable of improving compliance without external LGU subsidy. This provided computational evidence for the argument by \cite{Espino2025} that non-compliance in under-resourced areas was often driven by structural limitations rather than mere negligence.

\subsection{Heterogeneity of Barangay Operations and Behavior}
A critical finding from the data gathering phase (Appendix B) was that the Municipality of Bacolod could not be modeled as a monolithic entity. The simulation environment was characterized by profound heterogeneity across its seven constituent barangays, creating a non-convex optimization landscape for the Deep Reinforcement Learning (DRL) agent. A ``one-size-fits-all'' policy was mathematically destined to fail because the responsiveness of agents varied drastically depending on their local context. This supported the core argument of \cite{Zheng2022}, which posited that policy effectiveness was non-uniform and highly contingent on household heterogeneity, specifically income and education levels.

The economic landscape of the municipality was sharply divided between agrarian subsistence zones and emerging commercial centers, a disparity that fundamentally altered the marginal utility of money within the simulation’s utility function. According to the income profiles calibrated from Appendix B, Barangays Babalaya and Demologan represented the municipality's economic floor, with 80\% of their households classified as low-income earners. In the Agent-Based Model, this high concentration of poverty assigned a higher gamma value to these agents, meaning that monetary fines were disproportionately punitive, yet the opportunity cost of time (compliance effort) was critically high due to subsistence labor. This simulation outcome reinforced \cite{Badua2022}, who warned that financial penalties disproportionately affected low-income groups, rendering uniform punitive policies regressive and potentially unjust.

In stark contrast, Barangay Binuni exhibited a ``Wealthy Anomaly'' profile, comprising a mix of 50\% low-income, 30\% middle-income, and 20\% high-income households. This economic buffer reduced the friction of compliance, as wealthier households possessed the resources to manage waste segregation without sacrificing essential labor time. Barangay Ezperanza served as the emerging middle-class standard, with a balanced distribution that created a predictable response to economic incentives. This heterogeneity implied that a uniform fine of \textpeso 500 was mathematically devastating for a household in Mati (90\% low-income/farmers) but merely a nuisance for a household in Binuni, forcing the DRL agent to weigh the equity of enforcement across these disparate economic zones. This disparity highlighted the need for the ``robust optimization'' approach advocated by \cite{Gentile2022}, where policy had to be designed to withstand parameter uncertainty across diverse populations.

\begin{table}[htbp]
    \centering
    \caption{Calibrated Behavioral Weights and Legislative Allocation Strategies}
    \label{tab:behavioral_parameters}
    \resizebox{\textwidth}{!}{ 
    \begin{tabular}{l c c c c c c c c}
        \toprule
        & \multicolumn{4}{c}{{Agent Behavior Parameters}} & & \multicolumn{3}{c}{{Budget Allocation (\%)}} \\
        \cmidrule(lr){2-5} \cmidrule(lr){7-9}
        {Barangay} & {Attitude ($w_a$)} & {Norms ($w_{sn}$)} & {Effort ($C_e$)} & {Decay ($\gamma$)} & & {Enf} & {Inc} & {IEC} \\
        \midrule
        Binuni       & 0.65 & 0.80 & 0.64 & 0.02 & & 40 & 40 & 20 \\
        Ezperanza    & 0.40 & 0.70 & 0.55 & 0.03 & & 50 & 30 & 20 \\
        Babalaya     & 0.60 & 0.90 & 0.62 & 0.05 & & 90 & 5  & 5 \\
        Liangan East & 0.65 & 0.60 & 0.58 & 0.04 & & 25 & 65 & 10 \\
        Poblacion    & 0.55 & 0.20 & 0.48 & 0.03 & & 60 & 20 & 20 \\
        Demologan    & 0.60 & 0.60 & 0.62 & 0.05 & & 85 & 10 & 5 \\
        Mati         & 0.60 & 0.50 & 0.60 & 0.05 & & 30 & 50 & 20 \\
        \bottomrule
    \end{tabular}
    }
    \vspace{1ex}
    
    \raggedright
    \footnotesize
    \textit{Note: Agent behavior parameters correspond to Theory of Planned Behavior (TPB) weights: $w_a$ (Attitude), $w_{sn}$ (Social Norms), $C_e$ (Perceived Cost of Effort), and $\gamma$ (Compliance Decay Rate). Budget allocation percentages reflect the legislative constraints under the Status Quo scenario.}
\end{table}

Beyond economics, the distinct psychosocial profiles of each barangay determined how they responded to policy interventions. The simulation encoded these differences through TPB weights—Attitude, Social Norms, and Perceived Behavioral Control—and a decay factor representing habit retention. Barangay Poblacion presented the most significant behavioral challenge; despite its size, it exhibited an ``Urban Isolation'' profile with the lowest Social Norm weight (0.20) and the highest Cost of Effort (0.75). This indicated that social pressure from neighbors was ineffective in the high-density commercial center, and the friction to comply was maximal. This finding was consistent with \cite{Meng2018}, who demonstrated that recycling behavior was heavily dependent on neighborhood-level interactions; where such social cohesion was absent (as in high-density urban settings), Subjective Norms failed to drive compliance.

Conversely, Barangay Babalaya demonstrated the ``Tight-Knit Community'' paradox. It possessed the highest Social Norm weight (0.90) and Attitude (0.80), suggesting a population that deeply desired to comply due to strong community cohesion. However, this was offset by the rigid financial constraints mentioned previously. Furthermore, the model accounted for the persistence of education through variable decay rates. Barangay Mati, representing the agricultural sector, had the highest decay rate (0.10), implying that IEC campaigns were quickly forgotten as residents reverted to traditional practices. This necessitated the ``multi-pronged interventions'' recommended by \cite{Trushna2024}, validating that information campaigns had to be continuous rather than one-off events to combat behavioral decay.
